% 阿诺德·索末菲（综述）
% license CCBYSA3
% type Wiki

本文根据 CC-BY-SA 协议转载翻译自维基百科\href{https://en.wikipedia.org/wiki/Arnold_Sommerfeld}{相关文章}

阿诺德·约翰内斯·威廉·索末菲（Arnold Johannes Wilhelm Sommerfeld，德语发音：[ˈaʁnɔlt joˈhanəs ˈvɪlhɛlm ˈzɔmɐˌfɛlt]；1868年12月5日－1951年4月26日）是一位德国理论物理学家，他在原子物理与量子物理的发展中开创了先河，同时还培养并指导了新一代理论物理学的重要学生。他曾担任7位诺贝尔奖得主的博士导师与博士后导师，并指导过至少30位其他著名的物理学家和化学家。只有J. J. 汤姆孙的学生群体才能在成就上与之相提并论。

他引入了第二个量子数——方位量子数，以及第三个量子数——磁量子数。他还提出了精细结构常数，并开创了X射线波动理论。
\subsection{早年生活与教育}
索末菲尔德于1868年出生在一个在普鲁士有着深厚家族根基的家庭。他的母亲策齐莉·马蒂亚斯（Cäcilie Matthias，1839–1902）是波茨坦一位建筑商的女儿。他的父亲弗朗茨·索末菲尔德（Franz Sommerfeld，1820–1906）是一名医生，出身于柯尼斯堡的一个显赫家庭。1822年，阿诺德的祖父自内陆迁居到柯尼斯堡，进入普鲁士王国的宫廷邮政秘书处任职。索末菲尔德在其家族所属的普鲁士福音派新教教会中接受基督教洗礼。虽然他本人并不虔信宗教，但他终生未曾放弃自己的基督信仰。\(^\text{[3][4][5]}\)

索末菲尔德在故乡东普鲁士柯尼斯堡的阿尔贝蒂纳大学学习数学和自然科学。他的博士论文导师是数学家费迪南德·冯·林德曼，\(^\text{[5]}\)同时他还受益于数学家阿道夫·胡尔维茨和大卫·希尔伯特以及物理学家埃米尔·维歇特的授课。\(^\text{[7]}\)他曾参加学生兄弟会“德意志大学生社团”，并因此在决斗中在脸上留下了一道伤疤。\(^\text{[8]}\)1891年10月24日（时年22岁），他获得了博士学位。\(^\text{[9]}\)

获得博士学位后，索末菲尔德留在柯尼斯堡攻读教师资格。他于1892年通过了全国考试，随后开始了一年的兵役，服役于柯尼斯堡的预备役部队。他于1893年9月完成义务兵役，之后八年中，每年继续自愿服役八周。凭借翘起的胡须、健壮的体格、典型的普鲁士军人风度，以及脸上的剑伤，他给人留下了如同骠骑兵上校般的印象。\(^\text{[8]}\)
\subsection{职业生涯}
\subsubsection{哥廷根}
1893年10月，索末菲尔德前往哥廷根大学，该校当时是德国数学研究的中心。\(^\text{[10]}\)在那里，他通过一次幸运的私人关系，成为矿物学研究所西奥多·李比施（Theodor Liebisch）的助理——李比施曾是柯尼斯堡大学的教授，也是索末菲尔德家族的朋友。\(^\text{[11]}\)

1894年9月，索末菲尔德成为费利克斯·克莱因（Felix Klein）的助理，负责在克莱因的讲座期间进行全面的笔记整理，并为数学阅览室撰写讲义，同时还管理阅览室。\(^\text{[8]}\)在克莱因的指导下，索末菲尔德于1895年完成了他的《教授资格论文》，这使他得以在哥廷根大学成为一名私人讲师。\(^\text{[13]}\)

作为一名私人讲师，索末菲尔德讲授了广泛的数学及数理物理课程。他最早在哥廷根开设了偏微分方程课程，\(^\text{[8]}\)并在其整个教学生涯中不断发展，最终成为其教材系列《理论物理讲义》的第六卷，题为《物理中的偏微分方程》。\(^\text{[14]}\)
\begin{figure}[ht]
\centering
\includegraphics[width=6cm]{./figures/6f359ea6e05a3107.png}
\caption{1897年的索末菲尔德} \label{fig_ANOSf_1}
\end{figure}
1895年与1896年，克莱因关于旋转体的讲座促使克莱因与索末菲尔德合作撰写了四卷本著作《陀螺理论》——这一合作历时13年（1897–1910）。前两卷探讨理论，后两卷则涉及其在地球物理学、天文学和技术中的应用。\(^\text{[8]}\)索末菲尔德与克莱因的合作，使他在思维方式上更倾向于应用数学，并在讲授艺术上深受影响。\(^\text{[14]}\)

在哥廷根期间，索末菲尔德结识了约翰娜·黑普夫纳），她是哥廷根博物馆馆长\(^\text{[15]}\)恩斯特·黑普夫纳的女儿。1897年10月，索末菲尔德开始担任克劳斯塔尔-策勒费尔德矿业学院的数学教席，接替威廉·维恩。这一职位的收入足以使他最终与约翰娜结婚。\(^\text{[7][8][10]}\)

应克莱因的请求，索末菲尔德还担任了《数学科学百科全书》第五卷的主编。这是一项重大的工作，从1898年持续到1926年。\(^\text{[8][13]}\)
\subsubsection{亚琛}
1900年，在克莱因的推动下，索末菲尔德受聘为皇家亚琛工业大学（Königliche Technische Hochschule Aachen，后来的亚琛工业大学 RWTH Aachen University）应用力学教席的特别教授。在亚琛期间，他发展了流体力学理论，并对此保持了长期的研究兴趣。后来，在慕尼黑大学，索末菲尔德的学生路德维希·霍普夫和维尔纳·海森堡分别以这一课题撰写了他们的博士论文。\(^\text{[7][8][10][13]}\)

在亚琛任职期间，索末菲尔德对滑动轴承润滑理论的研究作出了重要贡献，因此被邓肯·道森列为23位“摩擦学先驱人物”之一。\(^\text{[16]}\)
\subsubsection{慕尼黑}
自1906年起，索末菲尔德出任慕尼黑大学物理学正教授，并担任新成立的理论物理研究所所长。这一职位是由时任慕尼黑物理研究所所长的威廉·伦琴（Wilhelm Röntgen）为他推荐的。\(^\text{[17]}\)在索末菲尔德看来，这一任命仿佛是被召唤进入一个“特权的行动领域”。\(^\text{[14]}\)

直到19世纪末、20世纪初，实验物理在德国学界中一直被认为地位更高。然而，20世纪初，慕尼黑的索末菲尔德与哥廷根大学的马克斯·玻恩等理论家，凭借其早期的数学训练，扭转了这种局面，使数理物理，即理论物理，成为推动学科发展的主导力量，而实验物理则用于验证或推动理论的进展。\(^\text{[18]}\)

在索末菲尔德门下获得博士学位后，沃尔夫冈·泡利、维尔纳·海森堡和瓦尔特·海特勒先后成为玻恩的助手，并在当时快速发展的量子力学领域作出了重要贡献。

在慕尼黑长达32年的教学生涯中，索末菲尔德既讲授普通课程和专业课程，也主持研讨班与学术讨论会。[14]普通课程内容包括：力学、可变形体力学、电动力学、光学、热力学与统计力学，以及物理中的偏微分方程。这些课程每周上课四小时，冬季学期13周，夏季学期11周，面向已经修读过伦琴（后来是威廉·维恩）所开设的实验物理课程的学生。此外，每周还有两个小时的课堂专门用于问题讨论。专业课程则紧扣研究热点，往往与索末菲尔德的研究兴趣相关；这些课程的内容后来常常以科学论文的形式发表。专业讲座的目标在于直面理论物理中的前沿问题，使索末菲尔德和学生能够对问题形成系统理解，而不论他们是否能真正解决这些前沿提出的难题。[19]在研讨班与学术讨论会期间，学生会被指定阅读当下的最新文献，并准备口头报告。[14]自1942年至1951年，索末菲尔德致力于整理自己的讲义，以便出版。[10] 这些讲义最终以六卷本《理论物理讲义》的形式问世。
